\chapter{Fundamentos e resultados centrais}
\section{Da pergunta ao delineamento}
O DIC de decomposição contém 24 recipientes independentes, oito em cada nível de umidade. A resposta é a porcentagem de perda de massa seca após um período comum. Cada recipiente recebe o tratamento independentemente; uma única câmara por tratamento confundiria umidade e câmara. O DBC utiliza outra base: 18 unidades em seis blocos completos, com uma unidade por tratamento em cada bloco.

No estudo de aves, há 36 localidades, 12 por ambiente, com quatro horas de esforço agregado por local. Shannon é uma variável sintética com logaritmo natural; topografia é o desvio-padrão da elevação em uma janela comum. As categorias de ambiente não foram casualizadas. Padronizar esforço não garante detectabilidade constante.

\section{ANOVA: a hipótese aparece na decomposição}
O modelo de um fator é $Y_{ij}=\mu+\tau_i+\varepsilon_{ij}$. A hipótese nula é $H_0:\mu_B=\mu_M=\mu_A$, contra pelo menos uma média diferente. Para o teste clássico, os erros são independentes, normais, de média zero e variância comum. Se $N$ é o número de unidades e $a$ o de tratamentos:
\begin{align}
SQ_T&=\sum_i\sum_j(y_{ij}-\bar y)^2,
&SQ_{\mathrm{trat}}&=\sum_i n_i(\bar y_i-\bar y)^2,\\
SQ_E&=\sum_i\sum_j(y_{ij}-\bar y_i)^2,
&SQ_T&=SQ_{\mathrm{trat}}+SQ_E,\\
F&=\frac{SQ_{\mathrm{trat}}/(a-1)}{SQ_E/(N-a)}.
\end{align}
A razão compara a variação entre médias com a variação residual. No DIC, $F_{2,21}=34{,}181$, com $p\approx2{,}49\times10^{-7}$; o valor crítico a 5\% é 3,467. O F global não identifica os pares distintos. O valor-p é condicional à nula e ao modelo; não é a probabilidade de a hipótese ser verdadeira.

\section{Tukey: diferenças com incerteza simultânea}
Tukey considera a família de comparações par a par e fornece intervalos simultâneos das diferenças \parencite{rtukey}. Para o desenho balanceado:
\begin{equation}
(\bar y_i-\bar y_j)\;\pm\;
q_{1-\alpha;a,\nu}\sqrt{QM_E/n}.
\end{equation}
O quantil $q$ vem da amplitude studentizada, $\nu$ é o grau de liberdade residual e $n$ o número de unidades por grupo. Esses intervalos não são os IC individuais de cada média.
\begin{table}[H]\centering\small
\caption{Diferenças de perda de massa, em pontos percentuais, no DIC simulado.}
\begin{tabular}{lrrr}\toprule Comparação & Diferença & IC95\% inferior & IC95\% superior\\\midrule
Média menos Baixa & 7,775 & 2,390 & 13,160\\
Alta menos Baixa & 17,624 & 12,239 & 23,009\\
Alta menos Média & 9,849 & 4,464 & 15,234\\\bottomrule
\end{tabular}\end{table}
Todos os intervalos excluem zero neste cenário. Isso não define, por si, a relevância ecológica nem permite medir processos que não foram observados. Diferenças entre porcentagens são expressas em pontos percentuais; aumentos relativos exigiriam uma referência adicional.

\section{Regressão e teste F parcial}
Na regressão simples, $H'_i=\beta_0+\beta_1T_i+\varepsilon_i$, em que $T$ é topografia em metros. O ajuste executado foi:
\begin{equation}\widehat H'=1{,}77292+0{,}02601T,\qquad R^2=0{,}3716.\end{equation}
O teste de $\beta_1=0$ usa $F_{1,34}=20{,}104=t^2$. Essa identidade ocorre porque se testa um único coeficiente angular sob o mesmo modelo. Não demonstra causalidade.

No modelo múltiplo, $T_{10}=(T-15)/10$ e $P_{100}=(P-2400)/100$. Conservado é a categoria de referência. A estimativa de topografia é 0,26867 por 10 m, com IC95\% de 0,18384 a 0,35350, ajustada por chuva e ambiente. O teste parcial compara modelos aninhados ajustados às mesmas observações:
\begin{equation}
F_{\mathrm{parcial}}=
\frac{(SQ_{E,R}-SQ_{E,C})/q}{SQ_{E,C}/(N-p_C)}.
\end{equation}
$R$ denota o modelo reduzido, $C$ o completo, $q$ o número de parâmetros adicionais e $p_C$ o total de parâmetros no modelo completo. Para topografia, $F_{1,31}=41{,}721$. Para ambiente como termo conjunto de dois graus de liberdade, $F_{2,31}=21{,}363$. Uma ANOVA sequencial de um único modelo não representa automaticamente a mesma hipótese parcial.

\section{Diagnóstico e sensibilidade}
Examine resíduos versus ajustados, Q--Q, escala-localização, alavancagem e distância de Cook. Valores de referência como $4/n$ para Cook e $2p/n$ para alavancagem orientam a investigação; não são regras de exclusão. Modelos adequados ao desenho continuam necessários mesmo diante de gráficos aparentemente satisfatórios \parencite{zuur2010}.

Welch é uma sensibilidade à heterogeneidade de variâncias, sem resolver dependência. A covariância HC3 mantém os coeficientes OLS e modifica sua incerteza; os intervalos robustos nesta atividade utilizam referência normal assintótica. Randomização no DIC testa uma nula aguda compatível com a alocação, diferente da simples igualdade de médias. A análise leave-one-out avalia sensibilidade e previsão interna; não é validação regional externa.

\begin{dica}
Na redação, combine desenho, unidade, estimativa, intervalo e limitação. ``Significativo'' não substitui a magnitude. Um resultado observacional ajustado continua exigindo cautela causal; um dado simulado demonstra o método, não um padrão real da região.
\end{dica}
